\chapter*{Webographie}
\addcontentsline{toc}{chapter}{Webographie}

\begin{enumerate}

  % ── Cadre législatif et institutions béninoises ───────────────

  \item Secrétariat Général du Gouvernement (SGG Bénin) — Journal Officiel~: Loi N° 2017-20 et Loi N° 2020-35 portant Code du Numérique en République du Bénin.\\
        \url{https://sgg.gouv.bj}\\
        \textit{Consulté le 18 août 2026.}

  \item APDP Bénin — Autorité de Protection des Données Personnelles~: guides, déclarations préalables et dispositions du Livre~V.\\
        \url{https://apdp.bj}\\
        \textit{Consulté le 25 août 2026.}

  \item ANSSI Bénin — Agence Nationale de la Sécurité des Systèmes d'Information~: cadre de confiance numérique et cybersécurité (Livre~IV).\\
        \url{https://anssi.bj}\\
        \textit{Consulté le 02 septembre 2026.}

  \item ARCEP Bénin — Autorité de Régulation des Communications Électroniques et de la Poste~: régulation des réseaux et des opérateurs (Livre~I).\\
        \url{https://arcep.bj}\\
        \textit{Consulté le 08 septembre 2026.}

  \item Ministère de la Justice et de la Législation du Bénin — Compétences judiciaires de la CRIET en matière de cybercriminalité (Livre~VI).\\
        \url{https://justice.gouv.bj}\\
        \textit{Consulté le 15 septembre 2026.}

  % ── Intelligence Artificielle, Embeddings & RAG ────────────────

  \item Ollama — Plateforme d'inférence locale souveraine et documentation d'exécution du modèle \textit{gemma3:1b}.\\
        \url{https://ollama.com}\\
        \textit{Consulté le 20 septembre 2026.}

  \item Hugging Face — Dépôt officiel du modèle \textit{OrdalieTech/Solon-embeddings-base-0.1} spécialisé dans le droit francophone.\\
        \url{https://huggingface.co/OrdalieTech/Solon-embeddings-base-0.1}\\
        \textit{Consulté le 24 septembre 2026.}

  \item Sentence Transformers — Documentation officielle pour l'encodage dense et les calculs de similarités vectorielles.\\
        \url{https://sbert.net}\\
        \textit{Consulté le 28 septembre 2026.}

  \item Hugging Face — Documentation du cross-encoder \textit{antoinelouis/crossencoder-camembert-base-mmarcoFR} pour le ré-ordonnancement contextuel.\\
        \url{https://huggingface.co/antoinelouis/crossencoder-camembert-base-mmarcoFR}\\
        \textit{Consulté le 02 octobre 2026.}

  \item Scikit-learn — Documentation sur la vectorisation sparse TF-IDF (\textit{TfidfVectorizer}) et la similarité cosinus.\\
        \url{https://scikit-learn.org/stable}\\
        \textit{Consulté le 05 octobre 2026.}

  % ── Backend, Base de données & Sécurité ────────────────────────

  \item FastAPI — Documentation officielle du framework web asynchrone haute performance en Python.\\
        \url{https://fastapi.tiangolo.com}\\
        \textit{Consulté le 18 août 2026.}

  \item Pydantic — Documentation de la validation de schémas typés et de données en Python (v2).\\
        \url{https://docs.pydantic.dev}\\
        \textit{Consulté le 10 septembre 2026.}

  \item SQLAlchemy — Documentation de l'ORM relationnel Python pour la modélisation des entités transactionnelles.\\
        \url{https://docs.sqlalchemy.org}\\
        \textit{Consulté le 18 septembre 2026.}

  \item SQLite — Spécifications et documentation officielle du moteur de base de données relationnelle autonome.\\
        \url{https://www.sqlite.org/docs.html}\\
        \textit{Consulté le 22 septembre 2026.}

  \item Python-Jose \& Passlib — Documentation de la sécurisation des échanges par jetons JWT (RFC 7519) et hachage fort par \textit{bcrypt}.\\
        \url{https://python-jose.readthedocs.io}\\
        \textit{Consulté le 04 octobre 2026.}

  % ── Frontend web & Expérience utilisateur ──────────────────────

  \item Next.js — Documentation officielle du framework React 14 avec architecture App Router.\\
        \url{https://nextjs.org/docs}\\
        \textit{Consulté le 20 août 2026.}

  \item React — Documentation officielle de la bibliothèque de composants d'interfaces utilisateur (React 18).\\
        \url{https://react.dev}\\
        \textit{Consulté le 28 août 2026.}

  \item Tailwind CSS — Documentation officielle du framework CSS utilitaire pour le design responsive et le mode sombre.\\
        \url{https://tailwindcss.com/docs}\\
        \textit{Consulté le 06 septembre 2026.}

  \item Lucide React — Bibliothèque officielle d'icônes vectorielles légères pour interfaces graphiques.\\
        \url{https://lucide.dev}\\
        \textit{Consulté le 14 septembre 2026.}

  % ── Publications et ressources scientifiques en ligne ──────────

  \item Lewis, P. et al. (2020) — \textit{Retrieval-Augmented Generation for Knowledge-Intensive NLP Tasks}. ArXiv.\\
        \url{https://arxiv.org/abs/2005.11401}\\
        \textit{Consulté le 12 octobre 2026.}

  \item Gao, Y. et al. (2023) — \textit{Retrieval-Augmented Generation for Large Language Models: A Survey}. ArXiv.\\
        \url{https://arxiv.org/abs/2312.10997}\\
        \textit{Consulté le 16 octobre 2026.}

  \item Ordalie Tech — Travaux de recherche et méthodologie d'évaluation des embeddings appliqués au droit en langue française.\\
        \url{https://ordalie.tech}\\
        \textit{Consulté le 20 octobre 2026.}

  \item EGUEDJI, Mahoutondji Obed Moriace — Mémoire consulté pour la conception du système RAG, l'architecture d'intelligence artificielle et la rédaction du présent mémoire.\\
        \textit{Consulté en septembre et en octobre 2026.}

\end{enumerate}